\subsection{One Solution, No Solution, or Infinitely Many}

\begin{examplebox}[One solution]
Two nonparallel lines meet in one point. The system
\[
\begin{cases}x+y=5,\\2x-y=1\end{cases}
\]
has the unique solution $(2,3)$.
\end{examplebox}

\begin{examplebox}[No solution]
The system
\[
\begin{cases}x+y=1,\\x+y=4\end{cases}
\]
describes parallel distinct lines. The same expression cannot equal both constants, so the system is inconsistent.
\end{examplebox}

\begin{examplebox}[Infinitely many solutions]
The system
\[
\begin{cases}x+y=3,\\2x+2y=6\end{cases}
\]
describes the same line twice. Every point on $x+y=3$ is a solution.
\end{examplebox}

\begin{interpretationbox}[Independence and compatibility]
One solution means the conditions determine one common point. No solution means the conditions conflict. Infinitely many solutions mean at least one condition is redundant and freedom remains.
\end{interpretationbox}
